\documentclass[11pt]{article}
\usepackage[T1]{fontenc}
\usepackage[utf8]{inputenc}
\usepackage{lmodern}
\usepackage[margin=1in]{geometry}
\usepackage{amsmath,amssymb,amsthm}
\usepackage{microtype}
\usepackage[colorlinks=true,linkcolor=blue,citecolor=blue,urlcolor=blue]{hyperref}
\newtheorem{theorem}{Theorem}
\newtheorem{lemma}{Lemma}
\newcommand{\T}{\mathbb T}
\newcommand{\R}{\mathbb R}
\newcommand{\dd}{\,\mathrm d}
\newcommand{\mi}{m_{\mathrm{in}}}
\newcommand{\A}{\mathcal A}
\newcommand{\J}{\mathcal J}
\newcommand{\G}{\mathcal G}
\newcommand{\D}{D}
\title{A half-power vanishing-noise estimate for a\\
quadratic kinetic mean-field game}
\author{Kaiwen Jin\\
\small\href{mailto:yui.ui.siki@gmail.com}{yui.ui.siki@gmail.com}}
\date{}
\begin{document}
\maketitle

\begin{abstract}
We consider a variational kinetic mean-field game with quadratic acceleration
cost, local quadratic congestion, and quadratic terminal congestion on
$\T^d\times\R^d$. For a smooth, compactly supported probability datum, the
density minimizers at velocity diffusivity $\nu$ and at zero diffusivity
are at distance $O(\nu^{1/2})$ in spacetime $L^2$ and terminal $L^2$, in
every dimension. The proof compares the two minimum values using kinetic
Gaussian averaging and an entropy identity for the current flux.
A positive convolution commuting with kinetic transport justifies this
identity in the original weak admissible class, including the noncompact
velocity domain and possible vacuum.
\end{abstract}

\section{The variational problem and the estimate}

Let $d\geq1$, $T>0$, and
\[
 Q=\T^d\times\R^d,\qquad S=(0,T)\times Q,\qquad
 \D=\partial_t+v\cdot\nabla_x.
\]
We use the unit-period torus with volume-one Haar measure.
Fix a smooth, nonnegative probability density $\mi$ on $Q$, compactly
supported in velocity. For $\nu\in[0,1]$, consider the minimum of
\begin{equation}\label{eq:cost}
 \J(m,a)=\A(m,a)+\G(m),\qquad
 \A(m,a)=\int_S\frac{|a|^2}{2m},\qquad
 \G(m)=\frac12\int_Sm^2+\frac12\int_Qm(T)^2,
\end{equation}
under
\begin{equation}\label{eq:pde}
 \D m-\nu\Delta_vm+\operatorname{div}_va=0,\qquad m(0)=\mi.
\end{equation}
The acceleration flux $a$ takes values in $\R^d$.
At vacuum, $|a|^2/m$ is zero when $a=0$ and infinite otherwise.
The density $m(t)$ is a narrowly continuous curve of probability densities,
has finite second velocity moment at every time, and satisfies
\eqref{eq:pde} distributionally; the cost \eqref{eq:cost} is finite.
In particular, the endpoint densities in \eqref{eq:cost} are the actual
narrow traces. The weak equation, including these traces, reads
\begin{equation}\label{eq:weak}
 \int_Qm(T)\phi(T)-\int_Q\mi\phi(0)
 =\int_S\left[
 m(\partial_t\phi+v\cdot\nabla_x\phi+\nu\Delta_v\phi)
 +a\cdot\nabla_v\phi\right]
\end{equation}
for smooth tests periodic in $x$ and compactly supported in $v$.
We write $(m_\nu,a_\nu)$ for minimizing pairs in this class and
$E_\nu=\J(m_\nu,a_\nu)$. All conclusions below hold for any such pairs.

Vanishing-viscosity rates for kinetic mean-field games are raised by
Tang and Zhang \cite[Section~7, item~4(b)]{TZ}. First-order kinetic games
with local couplings have a variational theory on the noncompact velocity
domain \cite[Abstract]{GPM}. Here we treat the specific quadratic problem
\eqref{eq:cost}--\eqref{eq:pde}. The comparison uses the entropy and
current-velocity action identity appearing in the elliptic setting in
Benamou, Carlier, Di Marino, and Nenna \cite[eqs.~(2.9)--(2.10)]{BCDN}.
We derive its kinetic version directly, without assuming differentiability
or Fisher information of an optimizer.

\begin{theorem}\label{thm:main}
For every integer $d\geq1$, every $T>0$, and every smooth, nonnegative,
compactly supported probability density $\mi$ on $Q$, the minimizing
densities defined above satisfy
\begin{equation}\label{eq:main}
 \|m_\nu-m_0\|_{L^2(S)}^2+
 \|m_\nu(T)-m_0(T)\|_{L^2(Q)}^2\leq4L\nu
 \qquad(0<\nu\leq1).
\end{equation}
Here the finite constant $L$, independent of $\nu$, can be chosen as
\begin{equation}\label{eq:constants}
 \begin{split}
 B&=\frac{T+1}{2}\|\mi\|_{L^2(Q)}^2,\qquad
 M_{\mathrm{in}}=\int_Q|v|^2\mi,\\
 K&=e^T(M_{\mathrm{in}}+2B+2dT),\\
 L&=|H_{\mathrm{in}}|+\frac d2\log(2\pi)+\frac{K+d}{2}+1,
 \qquad H_{\mathrm{in}}=\int_Q\mi\log\mi .
 \end{split}
\end{equation}
We use $0\log0=0$. In particular,
$\|m_\nu-m_0\|_{L^2(S)}\leq2\sqrt L\,\nu^{1/2}$.
\end{theorem}

The theorem is an estimate for the stipulated minimizing pairs.
No existence theorem for a classical value function is needed.
The exponent is independent of dimension; we do not assert its optimality.

Kinetic Gaussian averaging turns an inviscid pair into a noisy pair while
decreasing its cost. For the reverse comparison, the entropy identity
produces an inviscid competitor with the noisy minimizer's density and an
excess cost bounded by $L\nu$. We justify this construction by positive
convolution and restore the original datum before using minimality.
Strict convexity of the two density costs then bounds the squared $L^2$
distance by the same order.

\section{Moments and the noisy comparison}

Averaging a finite-action pair against a probability measure of translations
does not increase its action or either density-square cost. Indeed, weighted
Cauchy--Schwarz gives
\begin{equation}\label{eq:jensen}
 \frac{|\overline a|^2}{\overline m}
 \leq\overline{\frac{|a|^2}{m}},\qquad
 \overline m^{\,2}\leq\overline{m^2}.
\end{equation}
If $\overline m=0$, nonnegativity and finite action imply
$\overline a=0$, so the stated vacuum convention is preserved.
Also, finite action and total spacetime mass $T$ give
\begin{equation}\label{eq:flux-l1}
 \int_S|a|\leq\sqrt{2T\A(m,a)}.
\end{equation}

\begin{lemma}\label{lem:moment}
Every admissible pair with $0\leq\nu\leq1$ satisfies
\[
 \sup_{0\leq t\leq T}\int_Q|v|^2m(t)
 \leq e^T\left(M_{\mathrm{in}}+2\A(m,a)+2dT\right).
\]
\end{lemma}

\begin{proof}
Put $\psi_R(v)=|v|^2/(1+|v|^2/R^2)$ and
$Y_R(t)=\int_Q\psi_Rm(t)$. Direct differentiation gives
$|\nabla\psi_R|^2\leq4\psi_R$ and
\[
 \Delta\psi_R=
 \frac{2d}{(1+|v|^2/R^2)^2}
 -\frac{8|v|^2/R^2}{(1+|v|^2/R^2)^3}\leq2d.
\]
For fixed $R$, use $\psi_R$ in \eqref{eq:weak} after multiplication
by a velocity cutoff at radius $r$, and then let $r\to\infty$.
The extra flux term is bounded by
$CR^2r^{-1}\int_{\{|v|\geq r\}}|a|$, which vanishes by
\eqref{eq:flux-l1}. The diffusion cutoff terms are bounded by a constant
depending on $R$ times $r^{-1}$ or $r^{-2}$ times the integrated mass.
Endpoint terms pass by bounded convergence. Thus $Y_R$ is absolutely
continuous, with its representative fixed by narrow continuity, and
\[
 Y_R'\leq
 2\sqrt{Y_R}\left(\int_Q\frac{|a|^2}{m}\right)^{1/2}+2d
 \leq Y_R+\int_Q\frac{|a|^2}{m}+2d
 \quad\hbox{a.e.}
\]
Integrating after multiplication by $e^{-t}$ gives the claimed bound for
$Y_R(t)$. Finally $\psi_R\uparrow|v|^2$, so monotone convergence proves
the bound at every time. This argument does not require a previously
known time-integrable second moment.
\end{proof}

For $\nu>0$ and $t>0$, define the probability Gaussian on $\R^{2d}$
\begin{equation}\label{eq:kinetic-kernel}
 k_{\nu,t}(y,w)=
 \left(\frac{\sqrt3}{2\pi\nu t^2}\right)^d
 \exp\left(-\frac{|w|^2}{4\nu t}
           -\frac{3|y-tw/2|^2}{\nu t^3}\right).
\end{equation}
Its covariance blocks are
$\operatorname{Cov}(y)=2\nu t^3I_d/3$,
$\operatorname{Cov}(y,w)=\nu t^2I_d$, and
$\operatorname{Cov}(w)=2\nu tI_d$.
These follow by first integrating the conditional Gaussian in $y$,
of mean $tw/2$ and covariance $\nu t^3I_d/6$, then the Gaussian in $w$.
Direct differentiation shows
\[
 \partial_t k_{\nu,t}+w\cdot\nabla_yk_{\nu,t}
 -\nu\Delta_wk_{\nu,t}=0.
\]
Periodize $y$ on the torus, and denote convolution with this kernel by
$C_\nu(t)$, with $C_\nu(0)$ the identity. Integration by parts in the
shift variables yields
\begin{equation}\label{eq:noisy-operator}
 \D(C_\nu(t)f)=C_\nu(t)(\D f)+\nu\Delta_v(C_\nu(t)f).
\end{equation}
The identity holds distributionally as well: apply the adjoint convolution
to a smooth compact-velocity test and use the kernel equation.
Cutting off this adjoint test is justified by its Gaussian tails,
Lemma~\ref{lem:moment}, and \eqref{eq:flux-l1}.
The covariance tends to zero at $t=0$, so the original initial trace
is preserved.

The free-transport density $g(t,x,v)=\mi(x-tv,v)$ has constant $L^2$
norm and zero flux. Equation~\eqref{eq:noisy-operator} makes
$(C_\nu(t)g,0)$ an admissible noisy competitor. Together with
\eqref{eq:jensen}, this proves
\begin{equation}\label{eq:energy-bounds}
 0\leq E_\nu\leq B\quad(0\leq\nu\leq1),\qquad E_\nu\leq E_0.
\end{equation}
For the second inequality, use
$(C_\nu(t)m_0,C_\nu(t)a_0)$ instead. The averaged density is a
narrowly continuous probability curve and has velocity moment
$\int_Q|v|^2m_0(t)+2\nu dt$.
Narrow continuity follows by testing against bounded continuous functions:
the original curve is uniformly tight on $[0,T]$, and the Gaussian
shift law varies continuously, including its point-mass limit at zero.
All terms of the cost, including the terminal term, decrease under
this averaging.
Lemma~\ref{lem:moment} and \eqref{eq:energy-bounds} give the uniform bound
\begin{equation}\label{eq:uniform-moment}
 \sup_{0\leq t\leq T}\int_Q|v|^2m_\nu(t)\leq K
 \qquad(0\leq\nu\leq1),
\end{equation}
where $B$ and $K$ are defined in \eqref{eq:constants}.

For later use, set $H(f)=\int_Q f\log f$ for probability densities of
finite absolute entropy. If $f\in L^2(Q)$ and
$M(f)=\int_Q|v|^2f<\infty$, this entropy is absolutely finite.
The positive part is bounded by $f^2$. With
$\gamma(x,v)=(2\pi)^{-d/2}e^{-|v|^2/2}$ and
$c_\gamma=d\log(2\pi)/2$, the negative part on $\{f\leq1\}$ obeys
\[
 -f\log f\leq\gamma/e+f(c_\gamma+|v|^2/2).
\]
Integrating $r\log r-r+1\geq0$ against $\gamma$, with $r=f/\gamma$,
also gives
\begin{equation}\label{eq:entropy-lower}
 H(f)\geq-c_\gamma-\tfrac12M(f).
\end{equation}

\section{The entropy comparison in the weak class}

Fix $0<\nu\leq1$. Let $\theta_\varepsilon$ be the centered Gaussian
on $\R^d$ with covariance $\varepsilon^2I_d$, and let
$\rho_\varepsilon$ be its periodization on $\T^d$. Define
\begin{equation}\label{eq:shear-convolution}
 R_\varepsilon f(t,x,v)=
 \int_{\T^d\times\R^d}
 \rho_\varepsilon(y)\theta_\varepsilon(w)
 f(t,x-y-tw,v-w)\dd y\dd w,
 \quad
 q_\varepsilon=R_\varepsilon m_\nu,\quad
 c_\varepsilon=R_\varepsilon a_\nu .
\end{equation}
The shear $tw$ makes this convolution commute with kinetic transport:
for fixed $(y,w)$,
\[
 \D[f(t,x-y-tw,v-w)]=(\D f)(t,x-y-tw,v-w).
\]
It also commutes with velocity derivatives. Hence, distributionally,
\begin{equation}\label{eq:smoothed-pde}
 \D q_\varepsilon-\nu\Delta_vq_\varepsilon
 +\operatorname{div}_vc_\varepsilon=0.
\end{equation}
The density is smooth and positive in $(x,v)$ and has mass one.
The translations have unit Jacobian, so
\begin{equation}\label{eq:smoothed-bounds}
 \A(q_\varepsilon,c_\varepsilon)\leq\A(m_\nu,a_\nu),\qquad
 \G(q_\varepsilon)\leq\G(m_\nu),\qquad
 M(q_\varepsilon(t))=M(m_\nu(t))+d\varepsilon^2 .
\end{equation}
Its initial density is $R_\varepsilon\mi$ at $t=0$.
We will recover $\mi$ by taking $\varepsilon\to0$ before using
inviscid minimality.

We first justify the entropy identity for this positive density.
In physical variables the convolution kernel is
\[
 \kappa_{\varepsilon,t}(x,v)
 =\rho_\varepsilon(x-tv)\theta_\varepsilon(v).
\]
Weighted Cauchy--Schwarz in the periodizing sum gives
$I(\rho_\varepsilon):=\int_{\T^d}
|\nabla\rho_\varepsilon|^2/\rho_\varepsilon\leq d/\varepsilon^2$.
The velocity score of $\kappa_{\varepsilon,t}$ is
$-t\nabla\log\rho_\varepsilon(x-tv)+\nabla\log\theta_\varepsilon(v)$.
The integrated cross term is zero since
$\int_{\T^d}\nabla\rho_\varepsilon=0$.
Applying the same weighted inequality to the convolution gradient gives
\begin{equation}\label{eq:fisher}
 F_\varepsilon:=\int_S\frac{|\nabla_vq_\varepsilon|^2}{q_\varepsilon}
 \leq\frac{d(T+T^3/3)}{\varepsilon^2}<\infty .
\end{equation}

For fixed $\varepsilon$, the kernel is bounded, so $q_\varepsilon$
is uniformly bounded above. By \eqref{eq:uniform-moment}, at least half
the mass of $m_\nu(t)$ lies in one fixed velocity ball at every time.
The positive periodic kernel $\rho_\varepsilon$ has a positive minimum.
Restricting the convolution to that ball therefore gives
\begin{equation}\label{eq:log}
 q_\varepsilon(t,x,v)\geq
 c_{\varepsilon,*}e^{-C_\varepsilon(1+|v|^2)},\qquad
 |\log q_\varepsilon(t,x,v)|\leq C_\varepsilon(1+|v|^2),
\end{equation}
with positive constants depending on $\varepsilon$.
All kernel derivatives are bounded on $[0,T]$.
On each compact velocity set, \eqref{eq:smoothed-pde} thus puts
$\partial_tq_\varepsilon$ in $L^1_tL^\infty_{x,v}$, using
\eqref{eq:flux-l1}. The continuous representative furnished by convolution
agrees with the weak representative, including at the endpoints.
Positivity and the one-dimensional $W^{1,1}$ chain rule
\cite[Section~8.2]{Brezis} now justify differentiating
$q_\varepsilon\log q_\varepsilon$ on such compact sets.

Let $\chi_R(v)$ equal one on $|v|\leq R$, vanish on $|v|\geq2R$,
and satisfy $|\nabla\chi_R|\leq C/R$, where $R\geq1$.
Multiply \eqref{eq:smoothed-pde} by
$\chi_R(1+\log q_\varepsilon)$ and integrate.
The $x$-transport integrates to zero by periodicity.
The velocity integration by parts produces the principal mixed and Fisher
terms and the two boundary terms
\[
 \int_S(1+\log q_\varepsilon)c_\varepsilon\cdot\nabla\chi_R,\qquad
 -\nu\int_S(1+\log q_\varepsilon)
                  \nabla_vq_\varepsilon\cdot\nabla\chi_R.
\]
On the annulus supporting these terms,
$(1+|v|^2)/R\leq C(R^{-1}+|v|)$.
Both terms tend to zero because
\[
 \begin{split}
 \int_S|v||c_\varepsilon|
 &\leq\sqrt{2\A(q_\varepsilon,c_\varepsilon)}
            \left(\int_S|v|^2q_\varepsilon\right)^{1/2},\\
 \int_S|v||\nabla_vq_\varepsilon|
 &\leq\sqrt{F_\varepsilon}
            \left(\int_S|v|^2q_\varepsilon\right)^{1/2}
 \end{split}
\]
are finite; the corresponding unweighted integrals are finite as well.
The principal mixed term is absolutely integrable by
Cauchy--Schwarz and \eqref{eq:fisher}. Endpoint entropies pass by absolute
integrability, as established above. We obtain
\begin{equation}\label{eq:entropy-identity}
 H(q_\varepsilon(T))-H(q_\varepsilon(0))
 =\int_S c_\varepsilon\cdot\nabla_v\log q_\varepsilon
 -\nu F_\varepsilon .
\end{equation}
Only the smoothed density has been used in this chain rule.

Define the current flux
$b_\varepsilon=c_\varepsilon-\nu\nabla_vq_\varepsilon$.
Then $\D q_\varepsilon+\operatorname{div}_vb_\varepsilon=0$.
Its action is finite by \eqref{eq:smoothed-bounds} and \eqref{eq:fisher}.
Expanding the square and using \eqref{eq:entropy-identity} gives
\begin{equation}\label{eq:current-action}
 \A(q_\varepsilon,b_\varepsilon)
 =\A(q_\varepsilon,c_\varepsilon)
 -\nu[H(q_\varepsilon(T))-H(q_\varepsilon(0))]
 -\frac{\nu^2}{2}F_\varepsilon .
\end{equation}
Initial convolution decreases entropy by convexity of $s\log s$.
The terminal entropy is bounded below by \eqref{eq:entropy-lower}
and \eqref{eq:smoothed-bounds}. For $0<\varepsilon\leq1$, these give
\[
 H(q_\varepsilon(0))-H(q_\varepsilon(T))
 \leq H_{\mathrm{in}}+c_\gamma+(K+d)/2\leq L.
\]
Keeping the negative Fisher term in \eqref{eq:current-action} until
this point, and then discarding it, yields
\begin{equation}\label{eq:excess}
 \J(q_\varepsilon,b_\varepsilon)\leq E_\nu+L\nu
 \qquad(0<\varepsilon\leq1),
\end{equation}
with the constant $L$ in \eqref{eq:constants}.
This estimate is uniform in smoothing, although \eqref{eq:fisher}
itself is not.

\section{Removing the convolution and proving the rate}

We remove smoothing with $\nu$ fixed. Translation continuity and the
Gaussian shift moments in \eqref{eq:shear-convolution} imply
\begin{equation}\label{eq:density-limits}
 \begin{aligned}
 q_\varepsilon&\longrightarrow m_\nu
       &&\text{in }L^1(S)\cap L^2(S),\\
 q_\varepsilon(T)&\longrightarrow m_\nu(T)
       &&\text{in }L^2(Q),\\
 q_\varepsilon(0)&\longrightarrow\mi
       &&\text{in }L^1(Q).
 \end{aligned}
\end{equation}
For spacetime convergence, apply translation continuity at almost every
time and dominated convergence to the contracted $L^p$ norms.
The bounded shear on $[0,T]$ makes the shift size tend to zero uniformly
in probability.

H\"older's inequality gives the flux bound
\begin{equation}\label{eq:flux-compactness}
 \int_S|b_\varepsilon|^{4/3}
 \leq
 \left(\int_S\frac{|b_\varepsilon|^2}{q_\varepsilon}\right)^{2/3}
 \left(\int_Sq_\varepsilon^2\right)^{1/3}.
\end{equation}
Both factors are uniformly bounded by \eqref{eq:excess} and
\eqref{eq:smoothed-bounds}. Choose a sequence
$\varepsilon_n\downarrow0$ realizing the lower action limit $A_*$.
Reflexivity and weak sequential compactness in $L^{4/3}(S;\R^d)$
\cite[Theorems~3.18 and~4.10]{Brezis} give, along a subsequence,
$b_{\varepsilon_n}\rightharpoonup b$.
These spaces are defined on the full noncompact domain.
For example, the velocity tails also satisfy
\[
 \int_{\{|v|>R\}}|b_\varepsilon|
 \leq\sqrt{2\A(q_\varepsilon,b_\varepsilon)}
           \sqrt{T(K+d)/R^2}\qquad(0<\varepsilon\leq1).
\]

We check that the limiting flux has finite action with precisely the
original vacuum convention.
For every compactly supported continuous vector field $z$ on $S$,
completing a square and passing to the limit give
\begin{equation}\label{eq:dual-action}
 \Phi(z)-\frac12\|z\|_{L^2(\mu)}^2\leq A_*,
 \qquad
 \Phi(z)=\int_Sz\cdot b,\qquad
 \dd\mu=m_\nu\dd t\dd x\dd v .
\end{equation}
Replacing $z$ by each real multiple of itself shows
$|\Phi(z)|\leq\sqrt{2A_*}\|z\|_{L^2(\mu)}$.
Continuous compactly supported fields are dense in this weighted space.
To see this without a lower bound on $m_\nu$, first truncate a weighted
$L^2$ field in size and support. Approximate the resulting bounded field
in Lebesgue $L^2$ by continuous compactly supported fields and clip the
approximants to a common size bound. On $\{m_\nu\leq R\}$ the weighted
error is at most $R$ times the Lebesgue squared error; the remaining
error is bounded by a constant times
$\int_{\{m_\nu>R\}}m_\nu$, which tends to zero.
Lebesgue density on the torus follows by approximation in a fundamental
cell and periodization; the Euclidean density result is
\cite[Theorem~4.12]{Brezis}.

Consequently $\Phi$ extends to the Hilbert space $L^2(\mu;\R^d)$.
The Riesz representation theorem \cite[Theorem~5.5]{Brezis}
provides $u$ in that space with norm at most $\sqrt{2A_*}$ and
$\Phi(z)=\int_Sz\cdot u\dd\mu$.
Both $b$ and $um_\nu$ are locally integrable; equality against compact
smooth tests identifies them almost everywhere. Hence $b=um_\nu$,
including $b=0$ on $\{m_\nu=0\}$, and
\begin{equation}\label{eq:action-lower}
 \A(m_\nu,b)\leq A_*=
 \liminf_n\A(q_{\varepsilon_n},b_{\varepsilon_n}).
\end{equation}
In particular, $b$ is integrable by \eqref{eq:flux-l1}.

Now pass to the inviscid weak equation for $(q_{\varepsilon_n},
b_{\varepsilon_n})$. The gradient of any smooth compact-velocity test
belongs to $L^4(S)$, so the flux pairing passes under the weak convergence
above, including for tests reaching a time endpoint.
All density and endpoint terms pass by \eqref{eq:density-limits}.
The result is \eqref{eq:weak} with $\nu=0$, initial density exactly $\mi$,
and density curve $m_\nu$. The narrow continuity, probability property,
and finite moments are those of this original curve.
Strong convergence of both square costs and \eqref{eq:action-lower}
therefore show that $(m_\nu,b)$ is an inviscid admissible pair with
\begin{equation}\label{eq:inviscid-cost}
 \J(m_\nu,b)\leq E_\nu+L\nu\leq E_0+L\nu .
\end{equation}
The construction also identifies
$b=a_\nu-\nu\nabla_vm_\nu$ distributionally. No differentiability of
$m_\nu$ was assumed to construct it.

\begin{proof}[Proof of Theorem~\ref{thm:main}]
The midpoint of $(m_0,a_0)$ and $(m_\nu,b)$ is admissible for zero
diffusivity: the weak equation is linear, the initial traces agree, and
the probability and moment conditions are preserved.
Perspective convexity and the exact square identity in both density costs
give
\[
 \begin{split}
 E_0
 &\leq \J\left(\frac{m_0+m_\nu}{2},\frac{a_0+b}{2}\right)\\
 &\leq\frac{E_0+\J(m_\nu,b)}{2}
 -\frac18\|m_\nu-m_0\|_{L^2(S)}^2
 -\frac18\|m_\nu(T)-m_0(T)\|_{L^2(Q)}^2 .
 \end{split}
\]
Insert \eqref{eq:inviscid-cost} and rearrange to obtain
\eqref{eq:main}. In particular, for every stipulated minimizing pair
and every $0<\nu\leq1$,
\[
 \|m_\nu-m_0\|_{L^2((0,T)\times\T^d\times\R^d)}
 \leq2\sqrt L\,\nu^{1/2}.
\]
Here $L$ is given by \eqref{eq:constants}: it depends only on $d$, $T$,
the initial $L^2$ norm, second velocity moment, and entropy.
All are finite for the stated datum, and none depends on $\nu$.
Thus the exponent $\eta(d)=1/2$ applies in every dimension.
\end{proof}

\paragraph{AI-use disclosure.}
OpenAI Codex was used to develop the argument, examine sources, and prepare
this manuscript. Two additional fresh Codex contexts supplied reader and
referee feedback on an earlier frozen draft, followed by revision with Codex.
These AI readings used the same vendor and a shared filesystem; they do not
constitute independent human review or formal proof verification.
The argument is presented in full for mathematical review.

\begin{thebibliography}{9}
\bibitem{TZ}
W.~Tang and Y.~P.~Zhang,
\emph{The convergence rate of vanishing viscosity approximations for
mean field games},
SIAM Journal on Mathematical Analysis \textbf{57} (2025), 3217--3254.
\href{https://doi.org/10.1137/24M1640008}{doi:10.1137/24M1640008}.

\bibitem{GPM}
M.~Griffin-Pickering and A.~R.~M\'esz\'aros,
\emph{A variational approach to first order kinetic mean field games
with local couplings},
\href{https://arxiv.org/abs/2112.03141v3}{arXiv:2112.03141v3},
11 July 2022.

\bibitem{BCDN}
J.-D.~Benamou, G.~Carlier, S.~Di Marino, and L.~Nenna,
\emph{An entropy minimization approach to second-order variational
mean-field games},
\href{https://arxiv.org/abs/1807.09078v2}{arXiv:1807.09078v2},
29 March 2019.

\bibitem{Brezis}
H.~Brezis,
\emph{Functional Analysis, Sobolev Spaces and Partial Differential Equations},
Universitext, Springer, 2011.
\href{https://doi.org/10.1007/978-0-387-70914-7}
{doi:10.1007/978-0-387-70914-7}.
\end{thebibliography}
\end{document}
